\documentclass[12pt]{article}

\newcommand{\monofont}{Consolas}
\newcommand{\mathfont}{Latin Modern Math}
%---DOCUMENT MARGINS---
\usepackage{geometry} % Required for adjusting page dimensions and margins
\geometry{
	paper=a4paper, % Paper size, change to letterpaper for US letter size
	top=4.25cm, % Top margin
	bottom=2.5cm, % Bottom margin
	left=2.5cm, % Left margin
	right=2.5cm, % Right margin
	headheight=3cm, % Header height
	%footskip=1.5cm, % Space from the bottom margin to the baseline of the footer
	headsep=0.8cm, % Space from the top margin to the baseline of the header
	%showframe, % Uncomment to show how the type block is set on the page
}
\usepackage{cmap}

\usepackage[T2A]{fontenc}
\usepackage[bulgarian]{babel}
\usepackage{fontspec}
\setmainfont{Times New Roman}
\setsansfont{Times New Roman}
\setmonofont{\monofont}
\usepackage[math-style = TeX]{unicode-math}
\setmathfont{\mathfont}

\usepackage[nobottomtitles*]{titlesec}
\usepackage[dvipsnames]{xcolor}
\titleformat
{\section} % command
{\normalfont\fontsize{14}{14}\sffamily\bfseries} % format
{} % label
{0pt} % sep
{\hspace{-0pt}{\color{cyan}\rule{0.7\textwidth}{\parskip}}\\*[2pt]} % before-code
\titlespacing{\section}{0pt}{0em}{0em}
\titleformat
{\subsection} % command
{\fontsize{14}{14}\bfseries} % format
{} % label
{0pt} % sep
{\hspace{-24pt}{\color{cyan}\rule{0.5\textwidth}{0.5\parskip}}\\*[5pt]} % before-code
\titlespacing{\subsection}{\parindent}{0em}{0em}

\newcommand{\formatdate}[3]{#1#2#3} % formats as YYYYMMDD
\newcommand{\addauthor}[4]{%
	\newcommand{\targetdate}{%
		\formatdate{#2}{#3}{#4}
	}
	\newcommand{\currentdate}{%
		\formatdate{\the\year}{\ifnum\month<10 0\fi\the\month}{\ifnum\day<10 0\fi\the\day}
	}
	\ifnum\currentdate>\targetdate
		{\noindent\hfill\bfseries{#1}}
	\fi
}

\setlength{\parskip}{0.5em}
\setlength{\parindent}{24pt}
\sloppy

\usepackage{fancyhdr}
\pagestyle{fancy}
\usepackage{background}
\backgroundsetup{
	contents={\includegraphics[width=28cm, height=3.45cm]{./structure/header_background.png}}, scale=1, placement=top, opacity=0.8
}
\fancyhead[L]{
	\begin{minipage}{\textwidth}
		\includegraphics[width=2.8cm]{./structure/ruo.png}
		\vspace{0.35cm}
	\end{minipage}
}
\fancyhead[R]{
	\begin{minipage}{\textwidth}
		\flushright\includegraphics[width=2.8cm]{./structure/sap.png}
		\vspace{0.35cm}
	\end{minipage}
}
\newfontfamily\headerfont[
	ItalicFont = *i,
	BoldItalicFont = *bi
]{Arial}
\fancyhead[C]{
	\begin{minipage}{\textwidth}
		\color{white}\centering\small\textbf{
			\headerfont
			\textit{ОТКРИТО ПЪРВЕНСТВО НА СОФИЯ ПО ИНФОРМАТИКА}\\
			\textit{\contestDate}\\
			\textit{Група \group, \grade\ клас}
		}
	\end{minipage}
}
\renewcommand{\headrulewidth}{0pt}
\fancyheadoffset[L,R]{1.5cm}
\usepackage{lastpage}
\fancyfoot[C]{\thepage\ / \pageref*{LastPage}}

\raggedbottom

\usepackage{amsmath}
\usepackage{stmaryrd}

\usepackage{graphicx}
\graphicspath{{./images/}}
\usepackage[export]{adjustbox}
\usepackage{wrapfig}
\makeatletter
\patchcmd\WF@putfigmaybe{\lower\intextsep}{}{}{\fail}
\AddToHook{env/wrapfigure/begin}{\setlength{\intextsep}{0pt}}
\makeatother
\usepackage[inkscapearea=page, inkscapepath=./svg-inkscape]{svg}
\svgpath{{./images/}}

\usepackage{makecell}
\usepackage{tabularray}
\AtBeginEnvironment{table}{\vspace{-0.2cm}}
\AtEndEnvironment{table}{\vspace{-0.2cm}}
\usepackage{float}
\usepackage{placeins}
\usepackage{caption}
\captionsetup[table]{
	skip=0.25em, font=it,
	singlelinecheck=false, justification=justified, indention=-24pt,
	margin={24pt, 0pt}
}

\usepackage{enumitem}
\setlist{itemsep=-0.4em, leftmargin=\parindent, topsep=-\parskip}
\newcommand{\tabitem}{\indent~~\llap{\textbullet}~~}

\usepackage{hyperref}
\hypersetup{
	colorlinks=true,
	citecolor=blue,
	linkcolor=blue,
	urlcolor=cyan,
}

\usepackage{emoji}
\renewcommand{\bottomtitlespace}{3.5cm}

\newcommand{\contestDate}{3 ноември 2024 г.}
\newcommand{\group}{A}
\newcommand{\grade}{11-12}

\begin{document}

\section{Задача A3. ИЛИ/И}
\addauthor{Автор: Илиян Йорданов}{2024}{11}{03}

Като отявлен любител на побитовите операции Пешо намира следното предизвикателство в социалните мрежи. Даден е един масив и се пита дали може да се разделят елементите му на две непразни множества, така че \textit{побитовото или} на елементите от първото множество да е равно на \textit{побитовото и} на елементите на второто множество. Такъв масив Пешо нарича побитов. Той бързо се справя с тази закачка и даже измисля усложнение. Нека имаме масив $a$ с $N$ елемента - $a_1, a_2, \ldots, a_N$. Трябва да се отговори дали са побитови $Q$ негови подмасива, всеки зададен с два индекса $l$ и $r$ - $a_l, a_{l+1}, \ldots, a_r$. Помогнете на Пешо да реши тази задача, като напишете програма \textbf{\texttt{orand}}, която намира отговорите на въпросите.

\subsection{Вход}

От първия ред на стандартния вход се въвеждат целите числа $N$ и $Q$. От следващия ред се въвеждат $N$ числа - елементите на масива $a$. От последните $Q$ реда се въвеждат по две цели числа $l$ и $r$ - левият и десният край на подмасива за съответния въпрос.

\subsection{Изход}
За всяка заявка, по реда във входа, изведете "Yes"{} или "No"{} (без кавичките) в зависимост от това дали подмасивът е побитов.

\subsection{Ограничения}
\begin{itemize}
	\item $1\leq N, Q\leq 10^5$
	\item $0\leq a_i < 2^{30}$
\end{itemize}

\subsection{Подзадачи}
\begin{table}[H]
	\begin{tblr}{|Q[c,m]|Q[c,m]|Q[c,m]|Q[c,m]|Q[c,m]|}
		\hline
		\textbf{Подзадача} & \textbf{Точки} & \textbf{Необходими подзадачи} &
		$N$ & 
		\textbf{Допълнителни ограничения} \\
		\hline
		$1$ & $0$ & -- & $-$ & Примерът. \\ 
		\hline
		$2$ & $11$ & $1$ & $\leq 10^1$ & $-$ \\ 
		\hline
		$3$ & $30$ & $1-2$ & $\leq 10^2$ & $-$ \\
		\hline
		$4$ & $38$ & $1-3$ & $\leq 10^3$ & $-$ \\
		\hline
		$5$ & $10$ & $1-4$ & $\leq 2.10^4$ & $-$ \\
		\hline
		$6$ & $11$ & $1-5$ & $\leq 10^5$ & $-$ \\
		\hline
	\end{tblr}
	\caption*{Точките за дадена подзадача се получават само ако се преминат успешно всички тестове, предвидени за нея и необходимите подзадачи.}
\end{table}
\FloatBarrier

\subsection{Пример}
\begin{table}[H]
	\begin{tblr}{|l|l|X[j]|}
		\hline
		\textbf{Вход} & \textbf{Изход} & \textbf{Обяснение на примера} \\
		\hline
		\texttt{\makecell[lt]{6 4 \\
		13 15 13 5 1 7 \\ 
		1 5 \\ 
		2 5 \\ 
		3 6 \\
		4 6}}
		& 
		\texttt{\makecell[lt]{Yes \\
		No \\
		Yes \\
		No}}
		& 
		{Първата заявка се отнася за подмасива: $13, 15, 13, 5, 1$. Разделяме го на следните множества: $\{13, 5, 1\}$ и $\{13, 15\}$. \textit{Побитовото или} на първото множество е 13, колкото е \textit{побитовото и} на второто множество. \\
		Третата заявка се отнася за подмасива: $13, 5, 1, 7$. Разделяме го на следните множества: $\{5, 1\}$ и $\{13, 7\}$. \textit{Побитовото или} на първото множество е 5, колкото е \textit{побитовото и} на второто множество.} \\
		\hline
	\end{tblr}
\end{table}
\FloatBarrier

\end{document}